\begin{lemma}
A tripartite graph with $m$ edges has at most $(m/3)^{3/2}$ triangles.
\end{lemma}
\begin{proof}
Let the three partite classes be $A$, $B$, and $C$.  Since the graph is
tripartite, every edge joins two different classes.  Therefore every triangle
has one vertex in each of $A,B,C$: its three vertices have pairwise different
parts, and there are exactly three parts.

Write
\[
x=e(A,B),\qquad y=e(A,C),\qquad z=e(B,C).
\]
Then the total number of edges is $m=x+y+z$.  For a vertex $b\in B$, let
$a_b$ be the number of neighbors of $b$ in $A$, let $c_b$ be the number of
neighbors of $b$ in $C$, and let $t_b$ be the number of triangles whose
$B$-vertex is $b$.  Each such triangle is exactly an edge of the
$A$-$C$ bipartite subgraph with one endpoint in $N_A(b)$ and the other in
$N_C(b)$, so
\[
t_b\le a_b c_b
\qquad\text{and}\qquad
t_b\le y.
\]
For nonnegative numbers, $r\le p$ and $r\le q$ imply
$r^2\le pq$; applying this with $r=t_b$, $p=a_bc_b$, and $q=y$ gives
\[
t_b\le \sqrt{a_b c_b y}.
\]
Summing over $b\in B$ and using that each triangle has a unique vertex in
$B$, we get
\[
T=\sum_{b\in B}t_b
\le \sqrt y\sum_{b\in B}\sqrt{a_b c_b}.
\]
The Cauchy--Schwarz inequality gives
\[
\left(\sum_{b\in B}\sqrt{a_b c_b}\right)^2
\le
\left(\sum_{b\in B}a_b\right)
\left(\sum_{b\in B}c_b\right).
\]
This is the usual finite Cauchy--Schwarz inequality applied to the two
vectors $(\sqrt{a_b})_{b\in B}$ and $(\sqrt{c_b})_{b\in B}$; all entries are
nonnegative, so no absolute values change the expression.
Here $\sum_b a_b=x$, since both sides count the $A$-$B$ edges, and
$\sum_b c_b=z$, since both sides count the $B$-$C$ edges.  Hence
\[
T\le \sqrt y\,\sqrt{xz}=\sqrt{xyz}.
\]

Finally, by the arithmetic--geometric mean inequality,
\[
xyz\le \left(\frac{x+y+z}{3}\right)^3=\left(\frac m3\right)^3.
\]
The quantities $x,y,z$ are edge counts, hence nonnegative, so taking square
roots preserves the inequality.
Taking square roots of the nonnegative quantities yields
\[
T\le \sqrt{xyz}\le \left(\frac m3\right)^{3/2}.
\]
This is the claimed bound.
\end{proof}
